\documentclass[11pt,a4paper]{article}

\usepackage[utf8]{inputenc}
\usepackage[T1]{fontenc}
\usepackage{lmodern}
\usepackage[french]{babel}
\usepackage[strings]{underscore}
\usepackage[margin=2.2cm]{geometry}
\usepackage{tabularx}
\usepackage{longtable}
\usepackage{booktabs}
\usepackage{array}
\usepackage{colortbl}
\usepackage{xcolor}
\usepackage{enumitem}
\usepackage{amssymb}
\usepackage{listings}
\usepackage[colorlinks=true,linkcolor=blue!60!black,urlcolor=blue!60!black]{hyperref}

\hypersetup{
  pdftitle={PEKEGNO Management System -- Playbook d'execution developpement},
  pdfauthor={Equipe developpement Pekegno},
  pdfsubject={Specifications executables Backend / Frontend par sprint et par tache}
}

\setlength{\parskip}{0.35em}

% Couleurs
\definecolor{pekblue}{RGB}{54,65,245}
\definecolor{peklight}{RGB}{228,231,255}
\definecolor{pekgreen}{RGB}{21,150,110}
\definecolor{codebg}{RGB}{245,246,250}

% Blocs de code / DDL (sans accents pour compatibilite listings)
\lstset{
  basicstyle=\ttfamily\scriptsize,
  backgroundcolor=\color{codebg},
  frame=single,
  rulecolor=\color{pekgrayline},
  framerule=0.3pt,
  breaklines=true,
  columns=fullflexible,
  keepspaces=true,
  showstringspaces=false,
  aboveskip=0.6em,
  belowskip=0.6em,
}
\definecolor{pekgrayline}{RGB}{200,205,215}

\newcolumntype{Y}{>{\raggedright\arraybackslash}X}
\newcommand{\be}{\textsc{Be}}
\newcommand{\fe}{\textsc{Fe}}
\newcommand{\qa}{\textsc{Qa}}
\newcommand{\tache}[2]{\subsection{Tâche #1 --- #2}\label{task:#1}}
\newcommand{\critere}{$\square$ }

\begin{document}

%=====================================================
% Page de titre
%=====================================================
\begin{titlepage}
  \centering
  \vspace*{2.5cm}
  {\Huge\bfseries PEKEGNO\par}
  \vspace{0.3cm}
  {\LARGE MANAGEMENT SYSTEM\par}
  \vspace{1.8cm}
  {\Large\bfseries Playbook d'exécution\\[0.3cm] Spécifications Backend \& Frontend\\ par sprint et par tâche\par}
  \vspace{1.5cm}
  \begin{tabular}{rl}
    \textbf{Destinataires} & Agents de développement (opencode ou humain) \\[2mm]
    \textbf{Usage} & \og Exécute le sprint S3 \fg{} ou \og Exécute la tâche T2.1 \fg{} \\[2mm]
    \textbf{Référentiel} & Cahier des charges fonctionnel (Août 2026) \\[2mm]
    \textbf{Équipe} & Blackstar \& Le\_million (alternance BE/FE) \\
  \end{tabular}
  \vfill
\end{titlepage}

\tableofcontents
\newpage

%=====================================================
\section{Objet et mode d'emploi du document}
%=====================================================

Ce document est la \textbf{source de vérité technique} pour le développement restant de l'application PEKEGNO Management System. Il découpe le travail en \textbf{sprints} (S1 à S10) contenant des \textbf{tâches} identifiées (T0.1, T1.2, \ldots), chacune spécifiée avec ses modifications de base de données, ses endpoints, ses fichiers à créer ou modifier, et ses critères d'acceptation vérifiables.

\subsection{Grammaire de commandes}

Un agent de développement (ou un collègue utilisant opencode) peut donner l'une des instructions suivantes :

\begin{center}
\begin{tabularx}{\textwidth}{|l|Y|}
\hline
\rowcolor{peklight}\textbf{Instruction} & \textbf{Comportement attendu de l'agent} \\ \hline
\verb|Exécute le sprint S3| & Exécuter séquentiellement toutes les tâches du sprint 3, dans l'ordre de leur numéro, en respectant intégralement le protocole de la section~\ref{sec:protocole}. \\ \hline
\verb|Exécute la tâche T2.1| & Exécuter uniquement cette tâche, même protocole. \\ \hline
\verb|Reprend la tâche T7.5| & Vérifier d'abord l'état du dépôt (ce qui existe déjà), compléter ce qui manque, puis finir la tâche. \\ \hline
\verb|Vérifie le sprint S2| & Passer en revue tous les critères d'acceptation du sprint et produire un rapport de conformité sans rien modifier. \\ \hline
\end{tabularx}
\end{center}

\subsection{Règles d'or (à lire avant toute exécution)}

\begin{enumerate}[nosep]
  \item \textbf{Une tâche à la fois.} Ne jamais commencer une tâche tant que la précédente du même sprint n'est pas terminée et validée.
  \item \textbf{Ce document prime.} En cas de conflit entre le code existant et ce document, ce document gagne ; signaler l'écart dans le rapport final.
  \item \textbf{Zéro destruction.} Jamais de migration destructive sur des tables en production (toujours additive : ajout de colonnes, de tables). Renommer interdit ; déprécier autorisé.
  \item \textbf{Rien codé en dur.} Pays, agences, départements, catégories, seuils et statuts viennent de la base ou de \texttt{settings}.
  \item \textbf{Vente $\neq$ paiement.} Une vente possède N paiements ; interdiction absolue de créer des colonnes type \og tranche1 / tranche2 \fg{}.
  \item \textbf{Périmètre de tâche strict.} Ne pas retoucher des modules hors du périmètre de la tâche en cours, sauf point d'intégration explicitement décrit.
\end{enumerate}

%=====================================================
\section{Environnement et commandes}
%=====================================================

\subsection{Backend (Laravel 13, PHP 8.3, PostgreSQL)}

Le backend tourne via l'image Docker locale \verb|pekegno-php|. Depuis la racine du projet :

\begin{lstlisting}
# Installer les dependances PHP
sudo docker run --rm -v "$PWD":/app -w /app --network=host pekegno-php composer install --no-interaction

# Reinitialiser completement la base (migrations + seeders)
cd backend && sudo docker run --rm -v "$PWD":/app -w /app --network=host pekegno-php php artisan migrate:fresh --seed --force

# Regenerer la documentation Swagger apres toute modification d'API
sudo docker run --rm -v "$PWD":/app -w /app --network=host pekegno-php php artisan l5-swagger:generate --all

# Lancer les tests
sudo docker run --rm -v "$PWD":/app -w /app --network=host pekegno-php php artisan test

# Lister les routes / shell interactif
sudo docker run --rm -v "$PWD":/app -w /app --network=host pekegno-php php artisan route:list
sudo docker run --rm -it -v "$PWD":/app -w /app --network=host pekegno-php php artisan tinker

# Inspecter les tables PostgreSQL
sudo -u postgres psql -d pekegno -c "\dt"
sudo -u postgres psql -d pekegno -c "\d treasury_accounts"
\end{lstlisting}

\subsection{Frontend (React 19, Vite 6, Tailwind v4)}

\begin{lstlisting}
cd frontend
npm install          # premiere fois uniquement
npm run dev          # serveur de dev, proxy /api -> http://localhost:8000
npm run build        # DOIT passer sans erreur TypeScript (validation obligatoire)
\end{lstlisting}

Le build frontend (\verb|npm run build|) est le validateur minimal de chaque tâche frontend : il compile TypeScript en mode strict.

\subsection{Comptes de test}

Après \verb|migrate:fresh --seed| : connexion staff \verb|admin@pekegno.com| / \verb|password| (rôle super-admin, accès total). Dix utilisateurs de démonstration existent avec différents rôles (voir \verb|UserSeeder|).

%=====================================================
\section{Conventions obligatoires}
%=====================================================

Ces conventions reflètent le code existant. Les respecter à la lettre garantit l'homogénéité.

\subsection{Base de données}

\begin{itemize}[nosep]
  \item Clé primaire : \verb|$table->uuid('id')->primary()| ; modèles avec trait \verb|HasUuids|.
  \item Clés étrangères explicites : \verb|$table->foreign('x_id')->references('id')->on('table')->nullOnDelete()| (ou \verb|restrictOnDelete()| / \verb|cascadeOnDelete()| justifié).
  \item Montants : \verb|decimal('montant', 14, 2)|. Devise implicite FCFA, jamais stockée par ligne.
  \item Statuts métier : \verb|string| + contrainte CHECK PostgreSQL + constantes publiques dans le modèle (\verb|const STATUS_DRAFT = 'draft';|). Jamais de texte libre.
  \item Entités métier : \verb|timestamps()| + \verb|softDeletes()| + trio d'endpoints \verb|GET /x/trash|, \verb|POST /x/{id}/restore|, \verb|DELETE /x/{id}/force-delete|.
  \item Index sur chaque FK et sur les colonnes de filtrage fréquent.
  \item Transactions financières : suppression dure interdite. Annulation traçable (\verb|cancelled_at| + \verb|cancelled_by| + écriture miroir si trésorerie impactée).
  \item Commentaires français dans les migrations (\verb|->comment('...')|) pour toute ambiguïté.
\end{itemize}

Gabarit de migration canonique (respecter ce style) :

\begin{lstlisting}
<?php

use Illuminate\Database\Migrations\Migration;
use Illuminate\Database\Schema\Blueprint;
use Illuminate\Support\Facades\Schema;

return new class extends Migration
{
    public function up(): void
    {
        Schema::create('exemple_tables', function (Blueprint $table) {
            $table->uuid('id')->primary();
            $table->string('name', 150);
            $table->uuid('agency_id');
            $table->foreign('agency_id')->references('id')->on('agencies')->cascadeOnDelete();
            $table->string('status', 20)->default('draft');
            $table->timestamps();
            $table->softDeletes();

            $table->index('agency_id');
            $table->index('status');
        });

        // Contrainte CHECK specifique PostgreSQL si necessaire
        DB::statement("ALTER TABLE exemple_tables ADD CONSTRAINT exemple_tables_status_check
            CHECK (status IN ('draft','active','closed'))");
    }

    public function down(): void
    {
        Schema::dropIfExists('exemple_tables');
    }
};
\end{lstlisting}

\subsection{API backend}

\begin{itemize}[nosep]
  \item Toutes les routes staff vivent dans le groupe middleware \verb|['auth:sanctum','single.session','inactivity.logout','update.activity','portal:staff']| de \verb|routes/api.php|.
  \item Contrôle d'accès par middleware \verb|permission:<entite>.<action>|. Actions normalisées (9, voir \verb|PermissionSeeder|) : \verb|creer|, \verb|modifier|, \verb|supprimer|, \verb|consulter|, \verb|exporter|, \verb|imprimer|, \verb|valider|, \verb|encaisser|, \verb|annuler|. Toute nouvelle entité doit déclarer ses permissions dans \verb|PermissionSeeder| et les mapper aux rôles dans \verb|RoleSeeder|.
  \item Périmètres de visibilité via \verb|App\Services\ScopeService| : un utilisateur non-global ne voit que ses agences/départements assignés. Appliquer systématiquement dans les \verb|index| et filtres.
  \item Validation inline \verb|$request->validate([...])| avec messages français ; erreurs 422 au format Laravel.
  \item Logique réutilisable dans \verb|app/Services/*| (ex. PaymentService, CommissionService) ; contrôleurs minces.
  \item Journalisation des actions sensibles : trait \verb|App\Observers\Concerns\LogsActivity| ou \verb|ActivityLogger| (paiements, annulations, validations, changements de permissions).
  \item Annotations Swagger \verb|@OA\*| obligatoires sur chaque endpoint public, puis régénérer (\verb|l5-swagger:generate --all|).
  \item Listes paginées : \verb|per_page| (défaut 15), réponse au format paginate Laravel.
\end{itemize}

\subsection{Frontend}

\begin{itemize}[nosep]
  \item Un module API par domaine dans \verb|src/api/<domaine>.api.ts|, typé via \verb|src/types/<domaine>.ts|, consommant \verb|client| de \verb|src/api/client.ts|. Gabarit :
\end{itemize}

\begin{lstlisting}
import { client } from './client';
import type { Expense, ExpenseListParams } from '@/types/expense';
import type { PaginatedResponse } from '@/types/agency';

export const expensesApi = {
  async list(params: ExpenseListParams = {}): Promise<PaginatedResponse<Expense>> {
    const { data } = await client.get<PaginatedResponse<Expense>>('/expenses', { params });
    return data;
  },
  async create(payload: Partial<Expense>): Promise<Expense> {
    const { data } = await client.post<Expense>('/expenses', payload);
    return data;
  },
};
\end{lstlisting}

\begin{itemize}[nosep]
  \item Pages dans \verb|src/pages/<module>/|, importées en \emph{lazy} dans \verb|src/router/index.tsx| avec fallback \verb|SkeletonPage| (suivre le motif existant).
  \item Composants UI exclusivement depuis \verb|src/components/ui/| (Button, Input, Select, Modal, ConfirmDialog, Badge, Pagination, PeriodPicker, Autocomplete, Alert, Skeleton\ldots). Pas de nouvelle dépendance UI sans décision.
  \item Menus latéraux : \verb|src/components/layout/navItems.ts| (global) et fonctions par périmètre (\verb|getAgencyItems|, \verb|getDepartmentItems|\ldots). Chaque entrée : \verb|{ to, label: t('nav.cle'), icon: IconLucide, end: false }|.
  \item i18n : \textbf{les deux} fichiers \verb|src/i18n/locales/fr.ts| et \verb|en.ts| doivent contenir chaque nouvelle clé. Aucun texte littéral dans les composants.
  \item Montants : \verb|formatCurrency| (\verb|src/utils/number.ts|) ; dates : utilitaires de \verb|src/utils/date.ts|.
  \item Permissions côté client : utilitaires purs dans \verb|src/utils/*Permissions.ts| calqués sur les policies backend (indicatifs, le backend reste maître).
\end{itemize}

%=====================================================
\section{Architecture cible : navigation et contexte}\label{sec:navigation}
%=====================================================

\subsection{Modèle organisationnel corrigé}

\begin{lstlisting}
PEKEGNO GROUP (organization)
|-- country (ex: Cameroun)
|   |-- agency (ex: Douala)          <- structure physique
|   |   |-- department type=academy
|   |   |-- department type=agency   <- branche prestations
|   |   |-- department type=store
|   |   |-- department type=studio
\end{lstlisting}

Entrer dans une agence amène sur la page \textbf{Départements}. On y crée un département de l'un des quatre types. La notion actuelle d'\og agence academy/mixed \fg{} (\verb|agencies.type|) est \textbf{dépréciée} : ne plus ni l'écrire ni la lire. Le switcher \og Ligne Prestations / Ligne Formations \fg{} de \verb|AgencyLayout.tsx| est supprimé.

\subsection{Barre de contexte permanente (ContextBar)}

Présente en haut de toutes les vues authentifiées :

\begin{center}\ttfamily
PEKEGNO ~|~ Pays $\blacktriangledown$ ~|~ Agence $\blacktriangledown$ ~|~ Département $\blacktriangledown$ ~|~ Période $\blacktriangledown$ ~|~ Recherche ~|~ Notifications ~|~ Profil
\end{center}

\begin{itemize}[nosep]
  \item Alimentée par un arbre chargé une fois via \verb|GET /scope/context| (forme ci-dessous), mis en cache dans un provider React global \verb|OrgContext| (\verb|src/context/OrgContext.tsx|) persisté en localStorage.
  \item Sélection à n'importe quel niveau, y compris saut direct vers un autre pays : choisir \og Côte d'Ivoire $\rightarrow$ Abidjan \fg{} depuis Cameroun/Yaoundé/Academy doit rediriger proprement vers le nouveau contexte.
  \item La sélection pilote l'URL (routes paramétrées) et la période pilote tous les KPI.
\end{itemize}

\begin{lstlisting}
// GET /scope/context : forme cible
{
  "user": {"is_global": true, "assignments": [...]},
  "countries": [
    {
      "id": "...", "name": "Cameroun", "code": "CMR",
      "agencies": [
        {
          "id": "...", "name": "Yaounde", "code": "AG002",
          "departments": [
            {"id": "...", "name": "Academy Yaounde", "type": "academy"},
            {"id": "...", "name": "Store Yaounde", "type": "store"}
          ]
        }
      ]
    }
  ]
}
\end{lstlisting}

\subsection{Sidebars dynamiques par type de département}

\verb|getDepartmentItems(type)| retourne exactement ces menus :

\begin{itemize}[nosep]
  \item \textbf{academy} : Dashboard · Prospects · Apprenants · Inscriptions · Formations (Catalogue, Sessions, Planning, Formateurs) · Présences · Paiements · Créances · Certificats · Performances commerciales · Rapports
  \item \textbf{agency} : Dashboard · Prospects · Clients · Packages · Contrats · Prestations · Équipe client · Community Management · Publicité · Renouvellements · Paiements · Créances · Rapports clients · Performance
  \item \textbf{store} : Dashboard · Catalogue / Produits / Catégories · Stocks · Approvisionnements · Fournisseurs · Achats · Commandes · Ventes · Livraisons · Retours · Inventaires · Rapports
  \item \textbf{studio} : Dashboard · Prospects / Clients · Services · Devis · Projets · Briefs · Planning · Production · Révisions · Livraisons · Paiements · Rapports
\end{itemize}

Les pages existantes d'Academy (7 pages sous \verb|src/pages/academy/|) sont \textbf{réutilisées telles quelles} mais re-racines sous \verb|/departments/:departmentId/...|.

\subsection{Correspondance des routes (ancien $\rightarrow$ nouveau)}

\begin{tabularx}{\textwidth}{|Y|Y|}
\hline
\rowcolor{peklight}\textbf{Ancienne route} & \textbf{Nouvelle route} \\ \hline
/countries/:countryId/agencies/:agencyId (overview lignes) & /countries/:countryId/agencies/:agencyId (page Départements = landing) \\ \hline
\ldots/agencies/:agencyId/academy/* & /departments/:departmentId/(dashboard|prospects|learners|courses|sessions|trainers|presences|payments|receivables|certificates|reports) \\ \hline
\ldots/agencies/:agencyId/services & /departments/:departmentId/services (dept type agency) \\ \hline
/departments/:departmentId (actuel, 3 menus) & /departments/:departmentId/dashboard + menus dynamiques selon type \\ \hline
\end{tabularx}

Redirections permanentes des anciennes URLs maintenues pendant toute la durée du projet (composant de redirection dédié).

%=====================================================
\section{Modèle de données cible}\label{sec:modele}
%=====================================================

Spécification DDL complète des évolutions. Types abrégés : \verb|uuid| = UUID PK/FK, \verb|str(n)| = varchar(n), \verb|dec| = decimal(14,2), \verb|ts| = timestamps, \verb|sd| = softDeletes.

\subsection{M0 --- Organisation (Sprint 1)}

\begin{lstlisting}
ALTER departments
  + type           str(20) NOT NULL DEFAULT 'agency'
                   CHECK IN ('academy','agency','store','studio'), INDEX
  Backfill : departments heritant de agencies.type
             ('academy' -> academy, sinon 'agency').
DEPRECATED : agencies.type (ne plus lire/ecrire).

ALTER invoices     + department_id uuid NULL -> departments.id
ALTER orders       + department_id uuid NULL -> departments.id
  (contexte "branche" de chaque vente, requis par le noyau de ventes CDC)
\end{lstlisting}

\subsection{M1 --- Trésorerie (Sprint 2)}

\begin{lstlisting}
treasury_accounts
  id uuid PK
  agency_id      uuid NULL -> agencies.id nullOnDelete (NULL = compte groupe)
  name           str(100)
  type           str(20) CHECK IN ('cash','mobile_money','bank')
  provider       str(50) NULL  -- orange_money, mtn_momo, afriland, cca, other
  account_number str(50) NULL
  opening_balance dec DEFAULT 0
  currency_code  str(3) DEFAULT 'XAF'
  is_active      bool DEFAULT true
  ts + sd ; UNIQUE(agency_id, name) ; INDEX(type)

treasury_transactions
  id uuid PK
  treasury_account_id uuid -> treasury_accounts.id restrictOnDelete
  direction      str(10) CHECK IN ('in','out')
  amount         dec > 0
  source_type    str(50) NULL  -- morph: invoice_payment, expense, transfer, manual
  source_id      uuid NULL
  category       str(50) NULL  -- vente, approvisionnement, salaire, autre
  label          str(200)
  reference      str(100) NULL
  transacted_at  timestamp
  created_by     uuid -> users.id nullOnDelete
  ts ; INDEX(account_id, transacted_at) ; INDEX(source_type, source_id)

ALTER invoice_payments + treasury_account_id uuid NULL -> treasury_accounts.id NULL
  (backfill : compte cash de l'agence de la facture)
\end{lstlisting}

Invariant : le solde d'un compte n'est \textbf{jamas stocké} ; il vaut \verb|opening_balance + SUM(in) - SUM(out)|.

\subsection{M2 --- Dépenses (Sprint 3)}

\begin{lstlisting}
expenses
  id uuid PK
  number             str(30) UNIQUE  -- EXP-00001
  agency_id          uuid -> agencies.id cascadeOnDelete
  department_id      uuid NULL -> departments.id nullOnDelete
  category_id        uuid -> accounting_categories.id restrictOnDelete
                     (categories type='expense' uniquement, valide au save)
  amount             dec > 0
  expense_date       date
  status             str(20) DEFAULT 'draft'
                     CHECK IN ('draft','submitted','approved','rejected','paid','closed')
  requested_by       uuid -> users.id (demandeur)
  approved_by        uuid NULL -> users.id ; approved_at ts NULL
  rejected_by        uuid NULL -> users.id ; rejection_reason str(255) NULL
  paid_by            uuid NULL -> users.id ; paid_at ts NULL
  treasury_account_id uuid NULL -> treasury_accounts.id (compte debite au paiement)
  justification_path str(255) NULL  -- upload (photo/PDF)
  note               text NULL
  cancelled_at       ts NULL ; cancelled_by uuid NULL -> users.id
  ts + sd ; INDEX(status), INDEX(agency_id, expense_date)

Workflow : draft -> submitted -> approved -> paid -> closed
                       \-> rejected (raison obligatoire)
Transitions controlees dans ExpenseService ; passage a 'paid' =
  ecriture treasury_transactions(direction=out) + accounting_transactions
  (categorie depense) automatiques, atomique (DB::transaction).
\end{lstlisting}

\subsection{M3 --- Commissions versionnées (Sprint 3)}

\begin{lstlisting}
commission_rules
  id uuid PK
  rule_group_id   uuid  -- identite de la regle a travers ses versions
  version         int DEFAULT 1  -- incremente a chaque edition
  name            str(150)
  beneficiary_commercial_id uuid NULL -> commercials.id nullOnDelete
                  -- beneficiaire designe OU portee generique ci-dessous
  scope_country_id   uuid NULL -> countries.id
  scope_agency_id    uuid NULL -> agencies.id
  scope_department_id uuid NULL -> departments.id
  service_id         uuid NULL -> services.id  -- NULL = tout produit/service
  trigger_event   str(20) CHECK IN ('on_sale','on_payment','on_full_payment')
  formula_type    str(10) CHECK IN ('percent','fixed','tiered')
  percent_value   dec(5,2) NULL   -- si percent
  fixed_amount    dec NULL        -- si fixed
  tiers_json      jsonb NULL      -- si tiered : [{up_to, value, mode}]
  starts_on / ends_on date NULL
  is_active       bool DEFAULT true
  created_by      uuid -> users.id
  ts + sd ; INDEX(rule_group_id, version) UNIQUE(rule_group_id, version)

IMMUTABILITE : une regle publiee n'est jamais modifiee en place ;
l'edition cree une nouvelle ligne (version+1, meme rule_group_id)
et desactive la precedente.

commission_entries
  id uuid PK
  invoice_id        uuid -> invoices.id cascadeOnDelete
  invoice_payment_id uuid NULL -> invoice_payments.id nullOnDelete
  commission_rule_id uuid -> commission_rules.id restrictOnDelete
  rule_snapshot     jsonb  -- copie integrale de la regle appliquee
  beneficiary_commercial_id uuid -> commercials.id
  base_amount       dec
  amount            dec
  status            str(20) DEFAULT 'calculated'
                    CHECK IN ('calculated','validated','paid','cancelled')
  validated_by uuid NULL ; validated_at ts NULL
  paid_by      uuid NULL ; paid_at      ts NULL
  ts ; INDEX(invoice_id), INDEX(beneficiary..., status)

Plusieurs lignes possibles sur une meme vente/paiement.
CommissionService : a l'encaissement, evaluer TOUTES les regles actives
matchant le contexte ; conserver le fallback taux-fiche-commercial
(commission_payments existant) tant que aucune regle ne matche.
\end{lstlisting}

\subsection{M4 --- CRM complet (Sprint 4)}

\begin{lstlisting}
companies
  id uuid PK ; name str(150) ; industry str(100) NULL
  phone/email/address NULL ; city str(100) NULL ; country str(100) NULL
  website str(150) NULL ; ts + sd ; INDEX(name)

ALTER prospects + company_id uuid NULL -> companies.id nullOnDelete

opportunities
  id uuid PK
  title str(200)
  prospect_id   uuid NULL -> prospects.id cascadeOnDelete
  client_id     uuid NULL -> users.id nullOnDelete
  company_id    uuid NULL -> companies.id nullOnDelete
  agency_id     uuid -> agencies.id cascadeOnDelete
  department_id uuid NULL -> departments.id nullOnDelete
  commercial_id uuid -> commercials.id cascadeOnDelete
  stage         str(20) DEFAULT 'new' CHECK IN
                ('new','contacted','qualified','proposal','negotiation','won','lost')
  expected_amount dec NULL ; expected_close_date date NULL
  won_at / lost_at ts NULL
  ts + sd ; INDEX(stage), INDEX(commercial_id)

activities
  id uuid PK
  subject_type str(50) ; subject_id uuid  -- morph: Prospect, Opportunity, User(client)
  assigned_to  uuid -> users.id nullOnDelete
  created_by   uuid -> users.id nullOnDelete
  type    str(20) CHECK IN ('call','meeting','email','whatsapp','note','followup')
  title   str(200) ; notes text NULL
  due_at  ts NULL ; completed_at ts NULL
  outcome str(255) NULL
  ts ; INDEX(subject_type, subject_id), INDEX(assigned_to, due_at)

Job quotidien (routes/console.php) : rappels des activities echues non
completees -> notifications.
Dedoublonnage clients/prospects : recherche par telephone normalise
(chiffres seulement, suffixe 8 derniers chiffres).
\end{lstlisting}

\subsection{M5 --- Academy compléments (Sprint 5)}

\begin{lstlisting}
attendances
  id uuid PK
  training_session_id uuid -> training_sessions.id cascadeOnDelete
  enrollment_id       uuid -> enrollments.id cascadeOnDelete
  status str(10) DEFAULT 'present' CHECK IN ('present','absent','late','excused')
  recorded_by uuid -> users.id nullOnDelete ; recorded_at ts
  ts ; UNIQUE(training_session_id, enrollment_id)

certificates
  id uuid PK
  enrollment_id uuid -> enrollments.id cascadeOnDelete UNIQUE
  number  str(30) UNIQUE  -- CERT-00001
  issued_on date ; mention str(100) NULL
  status  str(10) DEFAULT 'issued' CHECK IN ('issued','revoked')
  revoked_reason str(255) NULL
  file_path str(255) NULL  -- PDF genere/imprimable
  created_by uuid -> users.id nullOnDelete
  ts + sd

ALTER enrollments + invoice_id uuid NULL -> invoices.id nullOnDelete
  (facture d'inscription ; solde apprenant = calcule depuis cette facture
   et ses paiements valides ; JAMAIS ressaisi)
\end{lstlisting}

\subsection{M6 --- Contrats Agency \& renouvellements (Sprint 5)}

\begin{lstlisting}
contracts
  id uuid PK
  number str(30) UNIQUE  -- CTR-00001
  client_id     uuid -> users.id restrictOnDelete (role client)
  company_id    uuid NULL -> companies.id nullOnDelete
  agency_id     uuid -> agencies.id cascadeOnDelete
  department_id uuid NULL -> departments.id nullOnDelete (dept type agency)
  pack_id       uuid NULL -> subscription_packs.id nullOnDelete
  start_date / end_date date
  billing_cycle str(20) DEFAULT 'monthly'
                CHECK IN ('one_shot','monthly','quarterly','yearly')
  amount dec
  status str(20) DEFAULT 'active' CHECK IN
         ('active','due_soon','expired','suspended','terminated')
  auto_renew bool DEFAULT false ; renewal_count int DEFAULT 0
  parent_contract_id uuid NULL -> contracts.id nullOnDelete (chaine historique)
  notes text NULL ; terminated_at ts NULL ; terminated_reason str(255) NULL
  ts + sd ; INDEX(client_id), INDEX(status, end_date)

contract_services
  contract_id uuid -> contracts.id cascadeOnDelete
  service_id  uuid -> services.id restrictOnDelete
  price dec NULL  -- surcharge du prix catalogue
  UNIQUE(contract_id, service_id)

ALTER accounting_categories + is_pass_through bool DEFAULT false
  -- true = budget publicitaire gere POUR le client : mouvement trace mais
  -- EXCLU du chiffre d'affaires (regle comptable CDC section 8.2)

Job quotidien 06:30 : end_date <= J+30 et status='active' -> 'due_soon'
  + notification ; end_date passee -> 'expired' + alerte.
Settings : renew_alert_days jsonb DEFAULT '[30,7]'
\end{lstlisting}

\subsection{M7 --- Store / Matootoo (Sprints 6--7)}

\begin{lstlisting}
suppliers
  id uuid PK ; code str(20) UNIQUE  -- FRN-00001
  name str(150) ; phone/email/address NULL ; city str(100) NULL
  notes text NULL ; agency_id uuid NULL -> agencies.id nullOnDelete
  is_active bool DEFAULT true ; ts + sd ; INDEX(name)

purchases                      -- Achats / Approvisionnements
  id uuid PK ; number str(30) UNIQUE  -- ACH-00001
  supplier_id  uuid -> suppliers.id restrictOnDelete
  agency_id    uuid -> agencies.id cascadeOnDelete
  department_id uuid NULL -> departments.id nullOnDelete
  status  str(20) DEFAULT 'draft' CHECK IN
          ('draft','ordered','partially_received','received','cancelled')
  order_date date ; expected_at date NULL ; received_at ts NULL
  subtotal dec ; discount dec DEFAULT 0 ; tax_rate dec(5,2) DEFAULT 0
  total_amount dec
  payment_status str(20) DEFAULT 'unpaid'
                 CHECK IN ('unpaid','partial','paid')
  treasury_account_id uuid NULL -> treasury_accounts.id  -- sortie fournisseur
  notes text NULL ; created_by uuid -> users.id
  ts + sd ; INDEX(supplier_id), INDEX(status, order_date)

purchase_lines
  id uuid PK ; purchase_id uuid -> purchases.id cascadeOnDelete
  product_id  uuid -> products.id restrictOnDelete
  qty_ordered dec(12,2) ; qty_received dec(12,2) DEFAULT 0
  unit_cost dec ; line_total dec
  ts

stock_levels
  id uuid PK
  product_id uuid -> products.id cascadeOnDelete
  agency_id uuid -> agencies.id cascadeOnDelete
  department_id uuid NULL -> departments.id nullOnDelete
  warehouse str(100) DEFAULT 'principal'
  qty_on_hand dec(12,2) DEFAULT 0
  reorder_point dec(12,2) DEFAULT 0
  ts ; UNIQUE(product_id, agency_id, department_id, warehouse)

stock_movements
  id uuid PK
  product_id uuid -> products.id restrictOnDelete
  agency_id uuid -> agencies.id restrictOnDelete
  department_id uuid NULL -> departments.id nullOnDelete
  type str(20) CHECK IN ('in','sale','out','return','loss','adjustment',
                         'transfer_in','transfer_out')
  quantity dec(12,2) > 0   -- toujours positif ; le sens vient du type
  unit_cost dec NULL ; reference str(100) NULL
  source_type str(50) NULL ; source_id uuid NULL  -- purchase_line, order_line...
  transfer_ref uuid NULL  -- couple transfert (out/in) partage cet uuid
  moved_at ts ; created_by uuid -> users.id
  ts ; INDEX(product_id, moved_at), INDEX(type), INDEX(transfer_ref)

StockService : seule porte d'ecriture du stock. Dans une DB::transaction :
  1) inserer stock_movements  2) recalculer qty_on_hand depuis les mouvements
  3) si qty_on_hand < reorder_point -> notification.
Types sortants : sale/out/loss/transfer_out. Entrants : in/return/transfer_in.
adjustment : delta signe via quantity positif + label.

ALTER order_lines + product_id uuid NULL -> products.id nullOnDelete
  (vente catalogue store ; line_type='catalog')

deliveries
  id uuid PK ; order_id uuid -> orders.id cascadeOnDelete
  status str(20) DEFAULT 'pending' CHECK IN
         ('pending','preparing','shipped','delivered','cancelled')
  carrier str(100) NULL ; tracking_number str(100) NULL
  delivered_at ts NULL ; received_by str(150) NULL
  notes text NULL ; ts ; INDEX(order_id, status)

returns                        -- Retours produits
  id uuid PK ; number str(30) UNIQUE  -- RET-00001
  order_id uuid NULL -> orders.id nullOnDelete
  invoice_id uuid NULL -> invoices.id nullOnDelete
  client_id uuid NULL -> users.id nullOnDelete
  product_id uuid -> products.id restrictOnDelete
  quantity dec(12,2) > 0 ; reason str(20) CHECK IN
           ('defective','wrong_item','client_remorse')
  restock bool DEFAULT true   -- reintegre le stock via StockService
  refund_amount dec DEFAULT 0 ; refund_method str(20) NULL
  treasury_account_id uuid NULL -> treasury_accounts.id  -- sortie remboursement
  status str(20) DEFAULT 'requested' CHECK IN
         ('requested','accepted','rejected','refunded','closed')
  processed_by uuid NULL -> users.id ; processed_at ts NULL
  ts + sd ; INDEX(status)

inventories
  id uuid PK ; number str(30) UNIQUE  -- INV-00001
  agency_id uuid -> agencies.id cascadeOnDelete
  department_id uuid NULL -> departments.id nullOnDelete
  status str(20) DEFAULT 'open' CHECK IN ('open','counting','closed')
  started_by uuid -> users.id ; closed_at ts NULL ; note text NULL
  ts + sd

inventory_lines
  id uuid PK ; inventory_id uuid -> inventories.id cascadeOnDelete
  product_id uuid -> products.id restrictOnDelete
  expected_qty dec(12,2)  -- gele a la creation de la session
  counted_qty dec(12,2) NULL
  unit_cost dec NULL
  adjustment_movement_id uuid NULL -> stock_movements.id nullOnDelete
  UNIQUE(inventory_id, product_id)

Cloture inventaire : ecart = counted_qty - expected_qty ; cree un
stock_movement type='adjustment' par ligne en ecart, valorise unit_cost.
\end{lstlisting}

\subsection{M8 --- Studio (Sprints 8--9)}

\begin{lstlisting}
quotes                          -- Devis
  id uuid PK ; number str(30) UNIQUE  -- DEV-00001
  client_id uuid NULL -> users.id nullOnDelete
  prospect_id uuid NULL -> prospects.id nullOnDelete
  company_id  uuid NULL -> companies.id nullOnDelete
  agency_id uuid -> agencies.id cascadeOnDelete
  department_id uuid NULL -> departments.id nullOnDelete
  status str(20) DEFAULT 'draft' CHECK IN
         ('draft','sent','accepted','refused','expired')
  valid_until date NULL ; sent_at ts NULL ; decided_at ts NULL
  subtotal dec ; discount dec DEFAULT 0 ; tax_rate dec(5,2) DEFAULT 0
  total_amount dec
  converted_project_id uuid NULL  -- pose a la conversion (pas de FK directe)
  created_by uuid -> users.id
  ts + sd ; INDEX(status)

quote_lines
  id uuid PK ; quote_id uuid -> quotes.id cascadeOnDelete
  label str(200) ; description text NULL
  unit_price dec ; quantity dec(12,2) DEFAULT 1 ; line_total dec
  position int DEFAULT 0 ; ts

projects
  id uuid PK ; number str(30) UNIQUE  -- PRJ-00001
  title str(200) ; description text NULL
  client_id uuid -> users.id restrictOnDelete
  company_id uuid NULL -> companies.id nullOnDelete
  quote_id uuid NULL -> quotes.id nullOnDelete
  agency_id uuid -> agencies.id cascadeOnDelete
  department_id uuid NULL -> departments.id nullOnDelete
  manager_id uuid NULL -> users.id nullOnDelete
  status str(20) DEFAULT 'planning' CHECK IN
         ('planning','in_production','in_revision','delivered','closed','cancelled')
  start_date date NULL ; due_date date NULL
  delivered_at ts NULL ; closed_at ts NULL
  budget_amount dec NULL
  invoice_deposit_id uuid NULL -> invoices.id nullOnDelete  -- acompte
  invoice_balance_id uuid NULL -> invoices.id nullOnDelete  -- solde
  ts + sd ; INDEX(status), INDEX(department_id)

project_briefs                  -- reunions de cadrage, notes riches
  id uuid PK ; project_id uuid -> projects.id cascadeOnDelete
  title str(200) ; meeting_date date
  participants text NULL        -- noms separes par virgule (v1)
  content_html longtext NULL    -- rendu Summernote (HTML assaini)
  created_by uuid -> users.id
  ts ; INDEX(project_id)

project_tasks                   -- planning / production
  id uuid PK ; project_id uuid -> projects.id cascadeOnDelete
  title str(200) ; description text NULL
  assignee_id uuid NULL -> users.id nullOnDelete
  start_date / due_date date NULL
  status str(20) DEFAULT 'todo' CHECK IN ('todo','in_progress','done','blocked')
  position int DEFAULT 0 ; completed_at ts NULL
  ts ; INDEX(project_id, status)

project_revisions
  id uuid PK ; project_id uuid -> projects.id cascadeOnDelete
  round int DEFAULT 1
  request_note text ; requested_by uuid -> users.id
  requested_at ts ; resolved_note text NULL ; resolved_at ts NULL
  status str(20) DEFAULT 'pending' CHECK IN ('pending','resolved')
  ts ; INDEX(project_id, status)

deliverables
  id uuid PK ; project_id uuid -> projects.id cascadeOnDelete
  title str(200) ; file_path str(255)
  revision_id uuid NULL -> project_revisions.id nullOnDelete
  version int DEFAULT 1
  uploaded_by uuid -> users.id
  delivered_at ts NULL
  ts ; INDEX(project_id)

Cycle CDC : Prospect -> Brief -> Devis -> Validation -> Acompte ->
Planification -> Production -> Revision -> Validation client ->
Livraison -> Solde -> Cloture.
\end{lstlisting}

\subsection{M9 --- Notifications centralisées (Sprint 9)}

\begin{lstlisting}
notifications (table generale, schema natif Laravel ok)
  id uuid PK ; user_id uuid -> users.id cascadeOnDelete
  type str(50)   -- receivable_overdue, renewal_due, stock_low,
                 -- expense_submitted, expense_decided, activity_reminder
  title str(200) ; body text NULL
  entity_type str(50) NULL ; entity_id uuid NULL
  read_at ts NULL
  timestamps ; INDEX(user_id, read_at)

NotificationService::notify(users[], type, title, body, entity)
Declencheurs a cabler :
  - immediat : expense submitted -> validateurs du perimetre ;
    stock sous seuil (apres StockService) ;
  - job quotidien 07:00 : creances echues (factures impayees echues),
    renouvellements contrats/abonnements J-X, rappels activities.
\end{lstlisting}

\subsection{M10 --- KPI, exports, outillage migration (Sprint 10)}

\begin{lstlisting}
import_logs
  id uuid PK ; filename str(255) ; entity str(50)
  mapping_version str(50) ; rows_total int ; rows_imported int
  rows_skipped int ; errors jsonb ; dry_run bool DEFAULT true
  executed_by uuid -> users.id ; ts

Commande : php artisan pekegno:import-sheets {file} {entity} [--dry-run]
Mappings : config/storage json versionne (storage/app/import-maps/*.json)
Regles : jamais d'ecriture comptable certifiee sans --force ;
   anomalies journalisees dans errors, lignes ignorees listees.
\end{lstlisting}

%=====================================================
\section{Endpoints API cibles}\label{sec:endpoints}
%=====================================================

Tous sous préfixe \verb|/api|, groupe staff authentifié. \verb|CRUD| signifie index/store/show/update/destroy + trash/restore/force-delete quand l'entité a \verb|softDeletes|.

\begin{longtable}{|p{1.3cm}|p{4.6cm}|p{2.7cm}|p{6.8cm}|}
\caption{Endpoints nouveaux ou modifiés par sprint.} \\
\hline
\rowcolor{peklight}\textbf{Verbe} & \textbf{Route} & \textbf{Permission} & \textbf{Description} \\
\hline
\endfirsthead
\hline
\rowcolor{peklight}\textbf{Verbe} & \textbf{Route} & \textbf{Permission} & \textbf{Description} \\
\hline
\endhead
\multicolumn{4}{|l|}{\cellcolor{peklight}\textit{S1 --- Contexte}} \\ \hline
GET & /scope/context & (auth) & Arbre pays$\rightarrow$agences$\rightarrow$départements du périmètre utilisateur, enrichi des types. \\ \hline
PUT & /scope/context-selection & (auth) & Persiste la dernière sélection de contexte de l'utilisateur (colonne \texttt{settings} JSON côté user ou table settings clé \texttt{context\_selection:\{userId\}}). \\ \hline
\multicolumn{4}{|l|}{\cellcolor{peklight}\textit{S2 --- Trésorerie}} \\ \hline
CRUD & /treasury/accounts & tresories.* & Comptes de trésorerie (cash/MoMo/banque). \\ \hline
GET & /treasury/accounts/\{account\}/transactions & tresories.consulter & Mouvements paginés d'un compte, filtres période/type. \\ \hline
POST & /treasury/transfers & tresories.valider & Transfert entre deux comptes (2 mouvements liés). \\ \hline
GET & /treasury/summary & tresories.consulter & Soldes courants par compte + totaux période. \\ \hline
GET & /bilans/daily & bilans.consulter & Bilan journalier \textbf{généré} (ventes par branche, encaissements par compte, dépenses, solde théorique vs réel). Paramètre \texttt{date}. \\ \hline
\multicolumn{4}{|l|}{\cellcolor{peklight}\textit{S3 --- Dépenses \& commissions}} \\ \hline
CRUD & /expenses & depenses.* & Dépenses. \\ \hline
POST & /expenses/\{expense\}/submit & depenses.modifier & draft $\rightarrow$ submitted. \\ \hline
POST & /expenses/\{expense\}/approve & depenses.valider & submitted $\rightarrow$ approved. \\ \hline
POST & /expenses/\{expense\}/reject & depenses.valider & submitted $\rightarrow$ rejected (raison requise). \\ \hline
POST & /expenses/\{expense\}/pay & depenses.encaisser & approved $\rightarrow$ paid : sortie trésorerie + écriture comptable. \\ \hline
CRUD & /commission-rules & commissions.* & Règles versionnées (update = nouvelle version). \\ \hline
GET & /commissions/entries & commissions.consulter & Lignes générées, filtres période/statut/bénéficiaire. \\ \hline
POST & /commissions/entries/\{entry\}/validate & commissions.valider & calculated $\rightarrow$ validated. \\ \hline
POST & /commissions/entries/\{entry\}/pay & commissions.valider & validated $\rightarrow$ paid. \\ \hline
\multicolumn{4}{|l|}{\cellcolor{peklight}\textit{S4 --- CRM}} \\ \hline
CRUD & /companies & entreprises.* & Entreprises. \\ \hline
CRUD & /opportunities & opportunites.* & Opportunités ; PATCH stage inclus. \\ \hline
CRUD & /activities & activites.* & Activités polymorphes ; POST \ldots/complete. \\ \hline
GET & /crm/timeline\{?\}subject\_type,subject\_id & crm.consulter & Timeline relationnelle unifiée. \\ \hline
\multicolumn{4}{|l|}{\cellcolor{peklight}\textit{S5 --- Academy \& contrats}} \\ \hline
PUT & /training-sessions/\{session\}/attendances & presences.modifier & Sauvegarde bulk feuille de présence. \\ \hline
CRUD & /certificates & certificats.* & Émission/révocation ; GET PDF imprimable. \\ \hline
CRUD & /contracts & contrats.* & Contrats Agency + services liés. \\ \hline
POST & /contracts/\{contract\}/renew & contrats.creer & Renouvelle : nouveau contrat chaîné, ancien $\rightarrow$ renewed/expired. \\ \hline
POST & /contracts/\{contract\}/terminate & contrats.annuler & Résiliation motivée. \\ \hline
\multicolumn{4}{|l|}{\cellcolor{peklight}\textit{S6--S7 --- Store}} \\ \hline
CRUD & /suppliers & fournisseurs.* & Fournisseurs. \\ \hline
CRUD & /purchases & achats.* & Achats/approvisionnements. \\ \hline
POST & /purchases/\{purchase\}/receive & achats.valider & Réception (qty par ligne) $\rightarrow$ mouvements \texttt{in}. \\ \hline
POST & /purchases/\{purchase\}/pay & achats.encaisser & Paiement fournisseur $\rightarrow$ sortie trésorerie. \\ \hline
GET & /stocks/levels & stocks.consulter & Niveaux, filtres produit/agence/département/seuil. \\ \hline
GET & /stocks/movements & stocks.consulter & Journal des mouvements. \\ \hline
POST & /stocks/movements & stocks.creer & Mouvement manuel (ajustement, perte, transfert). \\ \hline
POST & /orders/\{order\}/deliveries & livraisons.creer & Crée la livraison d'une commande. \\ \hline
PATCH & /deliveries/\{delivery\} & livraisons.modifier & Statut livraison. \\ \hline
CRUD & /returns & retours.* & Retours ; accept/refund intégrés au workflow. \\ \hline
CRUD & /inventories & inventaires.* & Sessions d'inventaire + lignes. \\ \hline
POST & /inventories/\{inventory\}/close & inventaires.valider & Clôture : ajustements de stock générés. \\ \hline
GET & /reports/store & rapports.consulter & Rotation, marge, top produits. \\ \hline
\multicolumn{4}{|l|}{\cellcolor{peklight}\textit{S8--S9 --- Studio \& notifications}} \\ \hline
CRUD & /quotes & devis.* & Devis ; POST \ldots/send, /accept, /refuse. \\ \hline
POST & /quotes/\{quote\}/convert & devis.valider & Convertit en projet (statut planning). \\ \hline
CRUD & /projects & projets.* & Projets studio. \\ \hline
CRUD & /projects/\{project\}/briefs & briefs.* & Briefs (notes HTML). \\ \hline
CRUD & /projects/\{project\}/tasks & projets.* & Tâches/planning ; réordonnancement via position. \\ \hline
CRUD & /projects/\{project\}/revisions & revisions.* & Cycles de révision. \\ \hline
CRUD & /projects/\{project\}/deliverables & projets.* & Livrables (upload). \\ \hline
POST & /projects/\{project\}/close & projets.valider & Clôture (solde facturé exigé). \\ \hline
GET & /notifications & (auth) & Liste personnelle, filtre non-lues. \\ \hline
POST & /notifications/read-all & (auth) & Marque tout lu. \\ \hline
PATCH & /notifications/\{notification\} & (auth) & Marque lu/non lu. \\ \hline
\multicolumn{4}{|l|}{\cellcolor{peklight}\textit{S10 --- Command Center \& imports}} \\ \hline
GET & /command-center & kpi.consulter & Cartes KPI consolidées + comparaison période précédente. Filtres : pays/agence/département/période. \\ \hline
GET & /kpi/drill-down & kpi.consulter & Descente dimensionnelle jusqu'à la transaction. \\ \hline
POST & /imports/preview & kpi.creer & Analyse CSV vs mapping, rapport d'anomalies (dry-run). \\ \hline
POST & /imports/run & kpi.creer & Import définitif journalisé (--force équivalent). \\ \hline
\end{longtable}

Permissions nouvelles à déclarer dans \verb|PermissionSeeder| (entités) : \verb|tresories|, \verb|depenses|, \verb|commissions| (si absent), \verb|entreprises|, \verb|opportunites|, \verb|activites|, \verb|presences|, \verb|certificats|, \verb|contrats|, \verb|fournisseurs|, \verb|achats|, \verb|stocks|, \verb|livraisons|, \verb|retours|, \verb|inventaires|, \verb|devis|, \verb|projets|, \verb|briefs|, \verb|revisions|, \verb|notifications|, \verb|kpi| --- avec les 9 actions normalisées quand pertinent. Mapping rôles (\verb|RoleSeeder|) : \verb|direction-generale| = tout ; \verb|responsable-agence| / \verb|responsable-departement| = tout sur leur périmètre ; \verb|comptable| = tresories, depenses, commissions, bilans ; \verb|caissier| = tresories consulter+encaisser, depenses consulter+creer ; \verb|commercial| = crm/opportunites/activites/devis consulter+creer+modifier.

%=====================================================
\section{Planning et alternance}
%=====================================================

Deux développeurs, Blackstar (BS) et Le\_million (LM). Sprint = 2 semaines ; les rôles BE/FE permutent à chaque sprint :

\begin{tabularx}{\textwidth}{|c|c|c|c|Y|}
\hline
\rowcolor{peklight}\textbf{Sprint} & \textbf{Semaines} & \textbf{Backend} & \textbf{Frontend} & \textbf{Chantier} \\ \hline
S1 & S1--S2 & BS & LM & Navigation \& contexte (T0.x) \\ \hline
S2 & S3--S4 & LM & BS & Trésorerie \& paiements (T1.x) \\ \hline
S3 & S5--S6 & BS & LM & Dépenses + commissions (T2.x) \\ \hline
S4 & S7--S8 & LM & BS & CRM complet (T4.x) \\ \hline
S5 & S9--S10 & BS & LM & Academy + contrats (T5.x) \\ \hline
S6 & S11--S12 & LM & BS & Store 1/2 (T7.x) \\ \hline
S7 & S13--S14 & BS & LM & Store 2/2 (T7.x) \\ \hline
S8 & S15--S16 & LM & BS & Studio 1/2 (T8.x) \\ \hline
S9 & S17--S18 & BS & LM & Studio 2/2 + notifications (T8/T9.x) \\ \hline
S10 & S19--S20 & LM & BS & Command Center \& imports (T10.x) \\ \hline
\end{tabularx}

Un développeur exécute les tâches de son flux (BE ou FE) ; les tâches \qa{} sont binômes. Quand une instruction porte sur un sprint entier, l'agent exécute les tâches correspondant à son flux puis signale celles du binôme.

%=====================================================
\section{Protocole d'exécution d'une tâche}\label{sec:protocole}
%=====================================================

À appliquer pour CHAQUE tâche, sans exception :

\begin{enumerate}
  \item \textbf{Lire} la carte de tâche et les sections référencées (modèle §\ref{sec:modele}, endpoints §\ref{sec:endpoints}, navigation §\ref{sec:navigation}).
  \item \textbf{Vérifier l'état} : \verb|git status| propre, \verb|migrate:fresh --seed| passe, \verb|npm run build| passe. Si une dépendance de la carte n'existe pas encore (tâche d'un autre flux non réalisée), implémenter le minimum syndical côté API (contrats définis ici) et le noter dans le rapport.
  \item \textbf{Backend} (si tâche BE) : migration(s) $\rightarrow$ modèle(s) avec constantes de statuts et relations $\rightarrow$ Service si logique métier $\rightarrow$ contrôleur + routes + middleware permission $\rightarrow$ ScopeService sur les listes $\rightarrow$ LogsActivity sur actions sensibles $\rightarrow$ annotations Swagger $\rightarrow$ mise à jour PermissionSeeder/RoleSeeder $\rightarrow$ seeder de démonstration réaliste.
  \item \textbf{Validation backend} : \verb|migrate:fresh --seed --force| sans erreur, \verb|php artisan test| vert, \verb|l5-swagger:generate --all|, vérifier les nouveaux endpoints via \verb|route:list|.
  \item \textbf{Frontend} (si tâche FE) : types TS $\rightarrow$ module \verb|*.api.ts| $\rightarrow$ page(s) avec kit UI $\rightarrow$ routes lazy + Skeleton $\rightarrow$ entrées de menu (navItems ou getDepartmentItems) $\rightarrow$ i18n fr \textbf{et} en $\rightarrow$ permissions utilitaires si besoin $\rightarrow$ états loading/error/empty.
  \item \textbf{Validation frontend} : \verb|npm run build| vert (TypeScript strict). Smoke test manuel du parcours critique décrit dans les critères.
  \item \textbf{Commit} : un commit par tâche, message \verb|feat(scope): TX.Y resume| (scope = backend, frontend, ou fullstack).
  \item \textbf{Rapport de tâche} (sortie finale obligatoire) : liste des fichiers créés/modifiés, endpoints exposés, migrations créées, écarts éventuels avec cette spécification, points ouverts.
\end{enumerate}

\subsection{Checklist Definition of Done globale}

\critere Migrations additives uniquement ; \verb|migrate:fresh --seed --force| passe. \\
\critere Statuts contrôlés (CHECK + constantes), aucun statut libre. \\
\critere Permissions \verb|entite.action| câblées et mappées aux rôles. \\
\critere Périmètre utilisateur respecté sur toutes les listes (ScopeService). \\
\critere Actions sensibles journalisées (activity\_logs). \\
\critere Swagger régénéré. \\
\critere i18n FR/EN complète, aucun texte en dur. \\
\critere \verb|npm run build| et \verb|php artisan test| verts. \\
\critere Solde/trésorerie/commission : calculés depuis les transactions sources, jamais ressaisis. \\
\critere Aucune donnée (pays, agence, département, catégorie, seuil) codée en dur.

%=====================================================
\section{SPRINT 1 --- Recadrage organisationnel \& navigation contextuelle}
%=====================================================

Objectif : appliquer le modèle Agence $\rightarrow$ Départements typés et livrer la ContextBar. Références : §\ref{sec:navigation}, modèle M0 (§\ref{sec:modele}).

\tache{T0.1}{Départements typés (BE, BS, 3 j)}
\textbf{Prérequis} : aucun. \\
\textbf{Travail backend} :
\begin{enumerate}[nosep]
  \item Migration additive : colonne \verb|type| sur \verb|departments| (M0) + backfill dans la migration (\verb|UPDATE departments JOIN agencies| selon \verb|agencies.type|).
  \item Modèle \verb|Department| : constantes \verb|TYPE_ACADEMY='academy'| etc., \verb|TYPES| array, validation dans DepartmentController \verb|in:academy,agency,store,studio|.
  \item Déprécier \verb|agencies.type| : retirer toute lecture/écriture (modèle \verb|Agency::deriveType|, \verb|TYPES|, affichages), garder la colonne. Commentaire de dépréciation.
  \item \verb|DepartmentController@store/update| : champ type requis/validé ; réponse inclut \verb|type|.
  \item Seeder : \verb|AgencySeeder| crée 4 départements (un par type) par agence de démonstration.
  \item Tests feature : création/refus type invalide, backfill correct.
\end{enumerate}
\textbf{Critères d'acceptation} :
\critere \verb|\d departments| montre la colonne type + CHECK. \\
\critere POST /departments sans type $\Rightarrow$ 422 ; avec \verb|"store"| $\Rightarrow$ 201 avec type persisté. \\
\critere Aucune référence active à \verb|agencies.type| hors migration historique (grep).

\tache{T0.2}{Arbre de contexte enrichi (BE, BS, 3 j)}
\textbf{Travail backend} : étendre \verb|ScopeController@__invoke| pour retourner la forme JSON de §\ref{sec:navigation} (pays $\rightarrow$ agences $\rightarrow$ départements avec \verb|id,name,type|), filtrée par \verb|ScopeService| ; ajouter option \verb|?types=academy,store|. Eager load performant (\verb|with| imbriqué). Tests : un utilisateur assigné à une seule agence ne voit que son sous-arbre.
\textbf{Critères} :
\critere Réponse conforme à la forme JSON documentée. \\
\critere Temps de réponse $<$ 300 ms avec données de démo. \\
\critere super-admin voit les 2 pays ; commercial limité.

\tache{T0.6}{Seeders \& Swagger (BE, BS, 2 j)}
Données de démonstration cohérentes : 2 pays, 3 villes, 3 agences, chacune avec ses 4 départements typés nommés proprement (\og Academy Douala \fg{}\ldots), utilisateurs de démo assignés à des départements précis. Swagger régénéré.

\tache{T0.3}{OrgContext + ContextBar (FE, LM, 5 j)}
\textbf{Travail frontend} :
\begin{enumerate}[nosep]
  \item Nouveau module \verb|src/api/scope.api.ts| (getContext, saveSelection).
  \item Provider \verb|src/context/OrgContext.tsx| : état \verb|{tree, selection:{countryId,agencyId,departmentId,period}, loading}|, persistance localStorage (\verb|pekegno_context|), hooks \verb|useOrgContext()|.
  \item Composant \verb|src/components/layout/ContextBar.tsx| : 4 dropdowns en cascade (Pays/Agence/Département/Période via \verb|PeriodPicker|), recherche globale (redirige vers page correspondante), cloche notifications (placeholder badge 0, lien futur), profil (réutilise \verb|UserMenu|).
  \item Intégration : monter la barre dans \verb|AppLayout|, \verb|CountryLayout|, \verb|AgencyLayout|, \verb|DepartmentLayout| (remplacer les headers back existants). Sélection = navigation vers l'URL canonique du niveau choisi ; changement de pays réinitialise agence/département.
\end{enumerate}
\textbf{Critères} :
\critere Depuis CMR/Yaoundé/Academy, ouvrir \og Pays \fg{} et choisir Côte d'Ivoire $\Rightarrow$ redirection immédiate vers Abidjan (ou liste agences CI) sans repasser par le dashboard. \\
\critere Le contexte survit au refresh (localStorage + URL cohérents). \\
\critere i18n fr/en complète.

\tache{T0.4}{Suppression du switcher de lignes + landing Départements (FE, LM, 4 j)}
\begin{enumerate}[nosep]
  \item \verb|AgencyLayout.tsx| : supprimer \verb|AgencyLine|, \verb|navigateLine|, \verb|getAgencyItems/getAcademyItems| et le segmented control ; la sidebar agence devient : Overview · Départements (défaut) · Équipes · Promotions · Abonnements · Services transverses · Settings.
  \item Landing agence = page Départements ; cartes de départements montrant le type (icônes : GraduationCap/Briefcase/ShoppingBag/Video), filtres par type.
  \item Modale de création : Select type (4 options traduites) + nom + description.
\end{enumerate}
\critere Plus aucune trace de \og Ligne Prestations/Formations \fg{}. \\
\critere Création d'un département de chaque type possible depuis l'UI.

\tache{T0.5}{Départements comme espaces de travail (FE, LM, 4 j + QA partagé)}
\begin{enumerate}[nosep]
  \item \verb|getDepartmentItems(type)| dans \verb|navItems.ts| retournant les menus §\ref{sec:navigation} ; items non encore implémentés pointent vers des routes placeholder propres (page \og À venir -- Sprint X \fg{} référençant la tâche).
  \item Ré-raciner Academy : routes \verb|/departments/:departmentId/{courses,sessions,trainers,learners,...}| réutilisant les composants existants (props \verb|departmentId| au lieu de \verb|agencyId| ; adapter les appels API qui filtrent déjà par \verb|agency_id| via query param).
  \item Redirections : \verb|/countries/:c/agencies/:a/academy/*| $\rightarrow$ département academy de l'agence (composant \verb|LegacyRedirect|).
  \item QA binôme : parcours Groupe $\rightarrow$ Pays $\rightarrow$ Agence $\rightarrow$ Département $\rightarrow$ cours Academy complet.
\end{enumerate}

%=====================================================
\section{SPRINT 2 --- Trésorerie \& bilan journalier}
%=====================================================

Références : modèle M1, endpoints §\ref{sec:endpoints}. Décisions requises en amont : liste des comptes réels (sinon seeder générique).

\tache{T1.1}{Comptes \& mouvements de trésorerie (BE, LM, 4 j)}
\begin{enumerate}[nosep]
  \item Migrations M1 (tables + colonnes ajoutées à \verb|invoices|, \verb|orders|, \verb|invoice_payments|) ; backfill \verb|invoice_payments.treasury_account_id| avec le compte cash de l'agence de la facture.
  \item Modèles + relations ; constantes \verb|TreasuryAccount::TYPE_*|, \verb|TreasuryTransaction::DIRECTION_*|.
  \item \verb|TreasuryService| : unique porte d'écriture \verb|recordMovement(account, direction, amount, source, ...)| en transaction ; méthode \verb|balanceFor(account)| ; \verb|transfer(from,to,amount)| = 2 mouvements liés.
  \item Branchement \verb|PaymentService| : chaque \verb|invoice_payments| enregistre le mouvement \verb|in| correspondant (montant, source=morph, référence numéro facture).
  \item Contrôleurs + routes + permissions \verb|tresories.*| ; Swagger ; tests (paiement $\Rightarrow$ solde juste ; annulation facture $\Rightarrow$ contrepassation).
\end{enumerate}
\critere Encaissement de 3 tranches sur une facture $\Rightarrow$ 3 mouvements \texttt{in}, solde exact. \\
\critere Impossible de payer une dépense depuis un compte inexistant/inactif (422).

\tache{T1.2}{Bilan journalier généré (BE, LM, 3 j)}
Refactorer \verb|BilanController| : \verb|GET /bilans/daily?date=&agency_id=| calcule depuis les sources : ventes par département type (via \verb|invoices.department_id|), encaissements par compte type, dépenses payées, solde théorique initial + in $-$ out ; compare à \verb|daily_balances.solde_final| saisi (écart + %). Conserver l'ancien endpoint compatible. Tests numériques exacts.

\tache{T1.4}{Encaissement avec compte (FE, BS, 3 j)}
Formulaire d'encaissement (InvoiceDetail + QuickSale) : Select compte de trésorerie (filtré agence, groupé par type avec icônes Banknote/Smartphone/Landmark), champ référence ; affichage colonne compte dans l'historique paiements. Types TS + i18n.

\tache{T1.5}{Page Trésorerie (FE, BS, 5 j)}
Menu FINANCE (global + pays + agence) : cartes soldes par compte, tableau flux paginé avec filtres (compte/type/période/Direction), modale transfert entre comptes, accès bilan jour. Permission \verb|tresories.consulter|.

\tache{QA-C1}{QA binôme (4 j)} Parcours CDC : vente 100~000 FCFA $\rightarrow$ 3 paiements $\rightarrow$ trésorerie $\rightarrow$ bilan cohérent ; transfert inter-comptes ; écarts affichés.

%=====================================================
\section{SPRINT 3 --- Dépenses \& moteur de commissions}
%=====================================================

Références : modèles M2, M3. Décision requise : formules réelles de prime (sinon règles de démo).

\tache{T2.1}{Module dépenses (BE, BS, 4 j)}
Tables M2 ; \verb|ExpenseService| (transitions contrôlées, refus motivé, paiement atomique : mouvement out + écriture comptable) ; contrôleurs/routes/permissions \verb|depenses.*| ; notification validateur (placeholder event, branchée au centre S9) ; upload justificatif via \verb|/uploads| ; tests workflow complet + refus de transition invalide.
\critere Transition \texttt{draft}$\rightarrow$\texttt{approved} directe impossible. \\
\critere Après \texttt{pay} : solde du compte baissé du montant exact, transaction comptable présente.

\tache{T2.2}{Règles de commissions versionnées (BE, BS, 4 j)}
Tables M3 ; extension \verb|CommissionService| : à l'encaissement, évaluer les règles actives matchant (périmètre, service, trigger) et créer les \verb|commission_entries| avec \verb|rule_snapshot| ; fallback existant préservé ; endpoints CRUD (update = nouvelle version) + entries + validate/pay ; tests multi-règles et paliers.
\critere Modifier une règle ne change PAS les entrées passées (snapshot). \\
\critere Deux règles matchant une même vente $\Rightarrow$ deux entrées.

\tache{T2.3}{UI Dépenses (FE, LM, 5 j)}
Pages : liste avec onglets par statut (badges colorés), formulaire demande (catégorie Select depuis \verb|accounting-categories| type expense, montant, date, compte visé, justificatif upload, note), file de validation (actions Valider/Rejeter avec raison), vue détail avec historique des statuts. Menu FINANCE/Dépenses, permissions alignées.

\tache{T2.4}{UI Commissions (FE, LM, 4 j)}
Page COMMISSIONS : onglet Règles (liste + modale création/édition : bénéficiaire, périmètres, formule percent/fixed/tiers avec éditeur de lignes de palier, dates, trigger ; avertissement \og édition = nouvelle version \fg{} + historique versions visible), onglet Commissions (entries paginées, filtres période/statut, actions Valider/Payer, aperçu du snapshot de règle appliquée).

\tache{QA-C2/C3}{QA binôme (1 j)} Parcours : règle 5\,\% Academy Douala $\rightarrow$ vente + encaissement $\rightarrow$ entry calculée $\rightarrow$ validation $\rightarrow$ paiement.

%=====================================================
\section{SPRINT 4 --- CRM complet}
%=====================================================

Référence : modèle M4. Réutilise \verb|prospects.api.ts| (déjà prêt).

\tache{T4.1}{Opportunités \& pipeline (BE, LM, 3 j)} Table opportunities + CRUD + PATCH stage (transitions libres tracées) ; conversion gagné $\Rightarrow$ possibilité créer client/facture liée (réutiliser convert prospect) ; scopes ; tests.

\tache{T4.2}{Activités \& relances (BE, LM, 3 j)} Table activities polymorphe + CRUD + \verb|POST /activities/{id}/complete| ; job quotidien de rappel (console.php 07:00) créant les notifications (service placeholder si M9 pas encore là : log + table prête).

\tache{T4.3}{Entreprises \& dédoublonnage (BE, LM, 2 j)} Table companies + CRUD + liaison prospect.company_id ; recherche clients/prospects normalisée téléphone (suffixe 8 chiffres) ; tests doublons.

\tache{T4.4}{Kanban prospects (FE, BS, 4 j)} Page PROSPECTS (manquante !) : kanban 7 colonnes (stages), drag \& drop HTML5 natif (aucune lib), carte = prospect/opportunité (nom, montant attendu, âge), filtres (commercial, agence, département, recherche), compteur par colonne ; menu CRM global + pays + agence.

\tache{T4.5}{Fiche prospect complète (FE, BS, 3 j)} Page détail : blocs Identité / Acquisition (source) / Intérêt (types cochables) / Responsable / Pipeline (stage + montant) / Activités (timeline + formulaire rapide appel/RDV/note/relance avec échéance) ; bouton Convertir en client (existante) ; i18n.

\tache{T4.6}{Timeline client (FE, BS, 2 j)} Onglet timeline sur ClientDetailPage : activités + factures + paiements + opportunités fusionnés chronologiquement (\verb|/crm/timeline|).

\tache{QA-C4}{QA binôme (1 j)} Parcours CDC prospect $\rightarrow$ qualification $\rightarrow$ proposition $\rightarrow$ gagné $\rightarrow$ client.

%=====================================================
\section{SPRINT 5 --- Academy compléments \& Contrats Agency}
%=====================================================

Références : modèles M5, M6.

\tache{T5.1}{Présences par session (BE, BS, 2 j)} Table attendances + \verb|PUT /training-sessions/{session}/attendances| (bulk upsert, valide enrollment appartient à la session) ; garde compat booléen \verb|enrollments.attendance| synchronisé (présent si dernière présence = present/late) ; tests.

\tache{T5.2}{Certificats (BE, BS, 2 j)} Table certificates ; émission si enrollment completed ; numéro CERT-##### ; révocation motivée ; route PDF imprimable simple (vue Blade print) ; permissions \verb|certificats.*| ; tests.

\tache{T6.1}{Contrats \& renouvellements (BE, BS, 4 j)} Tables M6 ; \verb|ContractService| : statuts calculés par job 06:30 (due\_soon à J-X depuis setting \verb|renew_alert_days|, expired après end\_date) ; renew (nouveau contrat chaîné parent\_contract\_id, renewal\_count++) ; terminate ; \verb|is_pass_through| sur catégories comptables + exclusion du CA dans StatsController (CA net du pass-through) ; tests renouvellement et exclusion CA.

\tache{T5.3}{UI Academy compléments (FE, LM, 5 j)} Feuille de présence session (tableau apprenants + 4 statuts radio, sauvegarde bulk, indicateur taux présence) ; onglet Certificats fiche apprenant (émis/revoqué, imprimable) ; bloc Finance fiche apprenant : échéancier = liste des paiements de la facture liée + solde calculé (déjà partiel, fiabiliser via \verb|enrollments.invoice_id|).

\tache{T6.2}{UI Contrats \& Renouvellements (FE, LM, 4 j)} Menu département type agency : Contrats (liste statuts colorés, modale création client+pack+services+prix+cycle+dates, détail avec services et historique renouvellements chaînés) ; Renouvellements (file due\_soon/expired avec actions Renouveler/Résilier, badges J-x).

\tache{QA-C5/C6}{QA binôme (1 j)} Parcours CDC : contrat $\rightarrow$ alerte J-7 simulée (job manuel tinker) $\rightarrow$ renouvellement $\rightarrow$ chaîne visible.

%=====================================================
\section{SPRINTS 6--7 --- Store / Matootoo}
%=====================================================

Référence : modèles M7. Tout le Store vit sous le département type \texttt{store}.

\tache{T7.1}{Fournisseurs \& achats (BE, LM, 4 j)} Tables suppliers/purchases/purchase\_lines ; receive (qty par ligne, statut partial/received, mouvements \texttt{in} via StockService, MAJ purchase\_price produit si fourni) ; pay fournisseur (sortie trésorerie, payment\_status) ; permissions ; tests réception partielle.

\tache{T7.2}{Stocks (BE, LM, 4 j)} Tables stock\_levels/stock\_movements + \verb|StockService| (transaction : mouvement $\rightarrow$ recalcul niveau $\rightarrow$ alerte seuil) ; endpoints levels/movements/movement manuel ; transfert inter-agences (paire transfer\_out/transfer\_in même transfer\_ref, 2 appels service atomiques) ; tests invariants (somme mouvements = niveau).

\tache{T7.3}{UI Catalogue \& Stocks (FE, BS, 5 j)} Sous dept store : Catalogue (produits existants API + catégories ; SKU, prix achat/vente, gestion stock on/off, dispo agence), Stocks (niveaux avec badge sous-seuil rouge, journal mouvements filtrable, modales ajustement/perte/transfert guidé 3 étapes).

\tache{T7.4}{UI Fournisseurs \& Achats (FE, BS, 4 j)} Fournisseurs CRUD ; Achats : liste statuts, création multi-lignes produits (Autocomplete produits), réception par ligne (qty reçue), paiement fournisseur, impression bon.

\tache{T7.5}{Commandes, livraisons, retours (BE, BS, 4 j)} order\_lines.product\_id ; commande store confirmée $\rightarrow$ réservation simple v1 ; deliveries CRUD/statuts ; returns workflow complet (accept $\rightarrow$ restock via StockService ; refund $\rightarrow$ sortie trésorerie) ; tests impacts stock/trésorerie.

\tache{T7.6}{Inventaires (BE, BS, 3 j)} Tables + open (gèle expected), counting (saisie counted), close (écarts $\rightarrow$ mouvements adjustment valorisés) ; permissions ; tests clôture avec écarts.

\tache{T7.7}{Rapports store (BE, BS, 1 j)} \verb|GET /reports/store| : rotation (qté vendue/stock moyen période), marge ((PV-PA)/PV pondérée), top produits, valeur du stock.

\tache{T7.8}{Ventes Store POS (FE, LM, 3 j)} Adapter QuickSale : onglets Services/Produits ; panier produits ; validation $\Rightarrow$ order + invoice + paiements + mouvements \texttt{sale} (backend orchestré dans OrderController confirm/invoice) ; stock temps réel dans le panier (badge rupture).

\tache{T7.9}{UI Commandes/Livraisons/Retours (FE, LM, 4 j)} Trois pages menu store : commandes (statuts, génération livraison), livraisons (kanban statuts simple + scan réf), retours (demande, traitement accept/reject, remboursement avec compte).

\tache{T7.10}{UI Inventaires \& rapports (FE, LM, 2 j)} Inventaires (sessions, saisie comptage mobile-friendly, écran écarts valorisés avant clôture) ; Rapports store (cartes rotation/marge/valeur stock + top produits).

\tache{QA-C7}{QA binôme (2+1 j)} Parcours CDC complet : achat 50 unités $\rightarrow$ vente 3 $\rightarrow$ retour 1 $\rightarrow$ inventaire $\rightarrow$ ajustement $\rightarrow$ rapports cohérents.

%=====================================================
\section{SPRINTS 8--9 --- Studio \& Notifications}
%=====================================================

Références : modèles M8, M9. Sémantique : projet = production concrète (ex. film publicitaire) ; devis = calcul du montant ; briefs = réunions avec notes Summernote.

\tache{T8.1}{Devis (BE, LM, 4 j)} Tables quotes/quote\_lines ; transitions send/accept/refuse datées ; expiration (valid\_until, job) ; convert $\rightarrow$ project planning (copie lignes en description) ; permissions \verb|devis.*| ; tests conversion.

\tache{T8.2}{Projets, briefs, planning (BE, LM, 4 j)} Tables projects/project\_briefs/project\_tasks ; réordonnancement tasks (position) ; statuts contrôlés ; sanitisation HTML briefs (\verb|purifier| ou strip\_tags liste blanche) ; permissions ; tests.

\tache{T8.3}{UI Devis \& Projets (FE, BS, 5 j)} Devis : liste pipeline coloré, éditeur lignes dynamique (label/prix/qté/remise/TVA, total live), aperçu imprimable, actions Send/Accept/Refuse/Convertir. Projets : cartes statuts (pipeline), détail avec onglets.

\tache{T8.4}{UI Briefs \& Planning (FE, BS, 4 j)} Briefs : liste réunions du projet + éditeur \textbf{Summernote} (installer \verb|react-summernote| + jQuery peer, décision validée car exigence métier ; fallback textarea HTML si build KO) ; participants, date ; rendu HTML consultable. Planning : tâches en vue calendrier mois simple (grille CSS native) + liste drag \& drop statuts.

\tache{T8.5}{Production, révisions, livraisons, clôture (BE, BS, 4 j)} Endpoints revisions (round auto-incrémenté par projet), deliverables (upload, version++), close (exige facture solde payée ou solde zéro ; passe delivered$\rightarrow$closed) ; acompte : à la conversion, générer facture acompte (50\,\% paramétrable settings) liée \verb|invoice_deposit_id| ; tests cycle complet CDC.

\tache{T8.6}{UI Production/Révisions/Livraisons (FE, LM, 5 j)} Onglets projet : Production (tâches en cours, avancement \%), Révisions (cycles demandeur/résolution), Livraisons (upload livrables versionnés, bouton livrer $\rightarrow$ statut delivered), Paiements (acompte/solde reliés factures, bouton encaisser réutilisant composant S2).

\tache{T9.1}{Centre de notifications (BE, BS, 4 j)} Table + \verb|NotificationService| ; brancher TOUS les déclencheurs M9 (dépense soumise, stock bas, créances échues, renouvellements, rappels activités, objectif) ; endpoints liste/read-all/patch ; job 07:00 consolidé ; rate-limit anti-spam (unique user+type+entité+jour).

\tache{T9.2}{Cloche \& centre (FE, LM, 4 j)} Cloche ContextBar (polling 60 s, badge count) ; panneau latéral notifications (non-lues en gras, clic $\Rightarrow$ navigation entité liée) ; page centre avec tout-lu et filtres ; i18n.

\tache{QA-C8/C9}{QA binôme (1+1 j)} Parcours CDC Studio complet prospect$\rightarrow$clôture + scénario notifications multi-déclencheurs.

%=====================================================
\section{SPRINT 10 --- Command Center, rapports \& migration}
%=====================================================

Référence : modèle M10. Décision requise : définitions KPI finales (sinon provisoires documentées).

\tache{T10.1}{API KPI consolidée (BE, LM, 4 j)} \verb|GET /command-center| : CA (net pass-through), encaissements, créances, dépenses, trésorerie par compte, performance (par pays/agence/dépt/commercial), alertes (top 10) ; chaque valeur = série comparable période précédente ; \verb|GET /kpi/drill-down?dimension=...| descente Groupe$\rightarrow$Pays$\rightarrow$Agence$\rightarrow$Dépt$\rightarrow$transaction (payload borné : 100 lignes max par niveau) ; cache 60 s.

\tache{T10.2}{Exports \& importeur Sheets (BE, LM, 4 j)} Etendre ExportController (command-center, treasury, expenses, commissions, contracts, store-report, studio-report) Excel/PDF/CSV ; commande \verb|pekegno:import-sheets {file} {entity} --dry-run|--\verb|force| avec mappings JSON versionnés, normalisation (dates, téléphones, montants), dédoublonnage par téléphone, import\_logs ; entités supportées v1 : clients, prospects, produits, formations.

\tache{T10.3}{Command Center (FE, BS, 5 j)} Redesign \verb|PEKEGNOGroupDashboard| : cartes KPI cliquables (delta \% vs période précédente, flèche tendance), drill-down au clic (modal/tableau niveau suivant jusqu'à la transaction avec lien facture), section Alertes actionable, tout piloté par ContextBar (pays/agence/dépt/période) ; graphiques : réutiliser le pattern CSS bars existant ou ajouter recharts (décision rapide à trancher, défaut : CSS).

\tache{T10.4}{Rapports transverses (FE, BS, 3 j)} Page RAPPORTS : filtres communs (période PeriodPicker, pays, agence, département, commercial) persistés dans l'URL ; onglets Direction/Commercial/Finance/Academy/Agency/Store/Studio consommant les endpoints existants + nouveaux ; boutons export.

\tache{T10.5}{Recette générale (QA, binôme, 2 j)} Exécuter les 7 parcours critiques du CDC de bout en bout sur données de démo ; corriger les bloquants trouvés ; checklist DoD §\ref{sec:protocole} complète ; rapport final.

%=====================================================
\section{Annexe --- Interdits et pièges connus}
%=====================================================

\begin{itemize}[nosep]
  \item Ne PAS réutiliser \verb|AdminDashboardPage| / \verb|DashboardPlaceholderPage| (orphelins legacy) : partir du composant existant group dashboard.
  \item \verb|CatalogSeeder| existe mais n'est pas appelé par \verb|DatabaseSeeder| : l'appeler si la tâche a besoin de catalogue de démo.
  \item Le frontend n'a AUCUN framework de test : ne pas en introduire dans une tâche métier (décision séparée).
  \item Pas de spatie/laravel-permission : RBAC maison (\verb|roles|, \verb|permissions|, \verb|role_permission|, \verb|users.role_id|).
  \item Portails Sanctum distincts staff/client (\verb|portal:staff|) : ne jamais exposer de routes business côté \verb|portal:client| sans décision.
  \item \verb|agencies.type| déprécié dès S1 ; ne plus s'y fier dans le moindre nouveau code.
  \item Les montants s'affichent en FCFA via \verb|formatCurrency| ; ne jamais formater à la main.
  \item Toute route nouvelle doit apparaître dans \verb|route:list| ET Swagger ; toute clé i18n doit exister dans fr.ts ET en.ts.
\end{itemize}

\end{document}
